\section{Related Work}
\label{sec:related}

\paragraph{Modality interference.}
Cai et al.~\cite{cai2025interference} study how signals from a
non-essential modality distort multimodal LLM decisions by applying discrete
perturbations to an irrelevant input and measuring argmax answer changes.
We instead degrade the informative modality along a continuous severity path
and relate the resulting distributional behaviour change to task utility.

\paragraph{Counterfactual grounding probes.}
Substituting blank or mismatched images to test visual dependence is
established practice. Zafar et al.~\cite{zafar2026grounding} compare real,
blank and shuffled images in medical VQA and report a visual reliance
score (accuracy with the real image minus accuracy with a shuffled one),
a blank-image accuracy drop, and an answer-change rate under shuffling;
they note that accuracy-based and answer-level metrics can disagree. Work
on over-reliance on visual context in VQA~\cite{gupta2022swapmix}
and on language priors in VLMs~\cite{goyal2017making,agrawal2018dont} is
closely related. Our swapped
condition at maximum severity is comparable to a shuffled-image probe. The
difference is that we trace a trajectory over severity, measure reliance
as divergence between full answer distributions rather than answer
changes, and locate the severity at which utility falls below the
no-image level.
OMD-Bench~\cite{nazi2026omd} replaces whole modality channels in
initially congruent video, audio and text inputs and evaluates
accuracy, calibrated abstention and confidence. It defines reliance as the
fractional accuracy drop when a modality is corrupted, and its authors
note that its binary corruption model may not transfer to graded
degradation. Our severity path addresses that regime directly, although
we restrict ourselves to one visual modality and a multiple-choice
format.

\paragraph{Quality-aware fusion.}
Moon et al.~\cite{moon2026quality} use positive controls to test whether
quality scores in fusion models influence decisions, separating outage from
degradation. We borrow the positive-control strategy in our synthetic
validation, but evaluate reliance directly on the degraded input rather than on
an auxiliary quality score.

\paragraph{Missing modalities and modality dropout.}
Robustness to absent modalities has been studied extensively
\cite{ma2022missing}; modality dropout as a training
technique~\cite{neverova2016moddrop} forces networks to learn robust
cross-modal features.
Our fusion study in Sect.~\ref{sec:fusion} trains small two-view models
with different modality-dropout rates and reports clean accuracy,
no-input accuracy and reliance; it is a proxy study and not a claim about
large multimodal systems.
